\input{configTD}
\usepackage{tikz}

% ── Poly à trous : commandes spécifiques ──────────────────────────
\newcommand{\atrou}[2][3cm]{%
  \ifthenelse{\boolean{showSolutions}}{%
    \par\vspace{0.5em}%
    \begin{mdframed}[backgroundcolor=ThemeLight, linewidth=0pt,
      skipabove=0pt, skipbelow=0pt,
      innertopmargin=8pt, innerbottommargin=8pt]%
    #2%
    \end{mdframed}%
    \par\vspace{0.5em}%
  }{%
    \par\vspace{0.5em}%
    \noindent\fbox{\begin{minipage}[c][#1]%
      {\dimexpr\textwidth-2\fboxsep-2\fboxrule\relax}%
      \color{white}\rule{0pt}{0pt}%
    \end{minipage}}%
    \par\vspace{0.5em}%
  }%
}

\newcommand{\val}[1]{%
  \ifthenelse{\boolean{showSolutions}}{#1}{\phantom{#1}}%
}

\newcommand{\R}{\mathbb{R}}

\title{\sffamily\bfseries Probabilités --- CM2\\[0.3em]
  \Large Variables aléatoires continues}
\author{}
\date{}

\begin{document}
\sffamily
\maketitle
\thispagestyle{empty}

% ====================================================================
\section*{1. Le problème : une mesure ne tombe jamais juste}
% ====================================================================

Jusque maintenant, nos variables aléatoires prenaient des valeurs
\textbf{entières} ou \textbf{isolées} : \newline 
\og combien de pièces
défectueuses ? \fg{} 
\og à quel rang apparaît le premier défaut ? \fg{}\newline
On pouvait alors lister les valeurs possibles $x_1,x_2,\dots$
et leur associer une probabilité $p_i=P(X=x_i)$.

\medskip

Dans la vraie vie, on mesure des grandeurs \textbf{continues} :
\textbf{physiques} :

\begin{center}
\begin{mdframed}[backgroundcolor=ThemeLight, linecolor=Theme,
  linewidth=1.5pt, innertopmargin=10pt, innerbottommargin=10pt]
\centering\large\itshape
Quelle est la résistance à 28 jours d'une éprouvette de béton ? \newline 
La longueur d'une pièce sortie d'atelier ? \newline 
Le temps d'attente avant la prochaine livraison ? \newline
\end{mdframed}
\end{center}

Ces grandeurs varient \textbf{continûment} : entre $34{,}1$ MPa
et $34{,}2$ MPa, il y a une infinité de valeurs possibles.

Faut-il alors écrire $P(X=34{,}17\ldots)$ ?
Et que vaut cette probabilité ?

\atrou[1.5cm]{%
Ces probabilités avec des égalités seront toutes nulles. 

On ne pourra calculer des probabilités que sur des intervalles.
}

\medskip

\textbf{L'objectif du cours.}

Adapter nos outils (espérance, variance, probabilités)
au cas où $X$ prend ses valeurs dans un \textbf{intervalle réel}.

Tout ce qu'on savait en discret a une \textbf{version continue} :\newline
les sommes deviennent des intégrales, \newline 
les probabilités ponctuelles
deviennent des \textbf{densités}.

% ====================================================================
\section*{2. De l'histogramme à la densité}
% ====================================================================

Imaginons qu'on mesure la longueur $X$ (en cm) de $1000$ pièces
sortant d'un atelier. On regroupe les mesures en classes et
on trace un \textbf{histogramme des fréquences}.

\medskip

On normalise l'histogramme pour que \emph{l'aire totale vaille~$1$}
(on divise la hauteur de chaque barre par la largeur de classe).

A quoi correspond l'aire d'une barre ?
\atrou[1.5cm]{%
L'aire d'une barre vaut alors la fréquence des mesures
tombées dans la classe.
}

\begin{center}
\begin{tikzpicture}[scale=0.9]
  % Histogramme grossier
  \draw[thick,->] (0,0) -- (9,0) node[right] {\small $x$};
  \draw[thick,->] (0,0) -- (0,3.2) node[above] {\small fréquence / largeur};
  \foreach \x/\h in {1/0.4, 2/0.9, 3/1.7, 4/2.5, 5/2.7, 6/2.1, 7/1.2, 8/0.5}{
    \draw[fill=ThemeLight, draw=Theme] (\x,0) rectangle (\x+1,\h);
  }
  \node at (4.5,-0.5) {\small largeur de classe : $1$ cm};
\end{tikzpicture}
\quad
\begin{tikzpicture}[scale=0.9]
  % Histogramme plus fin + courbe
  \draw[thick,->] (0,0) -- (9,0) node[right] {\small $x$};
  \draw[thick,->] (0,0) -- (0,3.2) node[above] {\small fréquence / largeur};
  \foreach \x/\h in {0.5/0.2, 1/0.35, 1.5/0.55, 2/0.8, 2.5/1.15, 3/1.55,
    3.5/1.95, 4/2.35, 4.5/2.65, 5/2.75, 5.5/2.55, 6/2.15, 6.5/1.7,
    7/1.25, 7.5/0.85, 8/0.5, 8.5/0.25}{
    \draw[fill=ThemeLight, draw=Theme] (\x,0) rectangle (\x+0.5,\h);
  }
  \draw[thick, Theme, domain=0.3:8.9, smooth, samples=80]
    plot (\x, {3*exp(-(\x-5)*(\x-5)/3.5)});
  \node at (4.5,-0.5) {\small on raffine les classes\dots};
\end{tikzpicture}
\end{center}

Quand on prend de plus en plus de mesures et des classes de
plus en plus fines, l'histogramme normalisé tend vers une
\textbf{courbe} $f$.

\begin{mdframed}[backgroundcolor=ThemeLight, linecolor=Theme,
  linewidth=1pt, innertopmargin=8pt, innerbottommargin=8pt, nobreak=true]
\textbf{Idée centrale.}
Pour une variable continue, ce qu'on associe à $X$ n'est plus
une \emph{liste de probabilités ponctuelles}, mais une \textbf{courbe}~$f$
(la densité) dont \textbf{l'aire sous la courbe} joue le rôle de probabilité :
\[
  P(a\le X\le b)=\text{(aire sous la courbe $f$ entre $a$ et $b$).}
\]
\end{mdframed}



Si $f(x_0)=2{,}7$, est-ce que \og la probabilité que $X$ vaille $x_0$ \fg{}
vaut $2{,}7$ ? \newline 
A quoi correspond cette valeur ? 

\atrou[1.5cm]{%
Évidemment non (une probabilité est $\le 1$).
Cette valeur n'est \textbf{pas} une probabilité,
c'est une \textbf{densité de probabilité}
(probabilité \emph{par unité de longueur}).

Cette valeur doit être comparée avec toutes les autres valeurs de la densité.
}

Pour obtenir une probabilité, il faut \textbf{intégrer}.

% ====================================================================
\section*{3. Densité : définition et premier usage}
% ====================================================================

\begin{mdframed}[backgroundcolor=ThemeLight, linecolor=Theme,
  linewidth=1pt, innertopmargin=8pt, innerbottommargin=8pt, nobreak=true]
\textbf{Définition.}
Une variable aléatoire $X$ est dite \textbf{continue}
si elle admet une \textbf{densité}, c'est-à-dire une fonction
$f:\R\to\R$ telle que :
\[
  f(x)\ge 0 \quad\text{pour tout } x\in\R,
  \qquad
  \int_{-\infty}^{+\infty} f(x)\,dx = 1,
\]
et pour tous réels $a\le b$,
\[
  \boxed{\,P(a\le X\le b)=\int_a^b f(x)\,dx.\,}
\]
\end{mdframed}

\textbf{Analogie avec le discret} (à garder en tête en permanence) :

\begin{center}
\renewcommand{\arraystretch}{1.5}
\begin{tabular}{|l|l|}
  \hline
  \textbf{Discret} & \textbf{Continu} \\
  \hline
  valeurs $x_i$ isolées & valeurs $x\in\R$ \\
  probabilités $p_i=P(X=x_i)$ & densité $f(x)$ \\
  $p_i\ge 0$ & $f(x)\ge 0$ \\
  $\displaystyle\sum_i p_i=1$ & $\displaystyle\int_{-\infty}^{+\infty} f(x)\,dx=1$ \\
  $\displaystyle P(X\in A)=\sum_{x_i\in A} p_i$
    & $\displaystyle P(X\in[a,b])=\int_a^b f(x)\,dx$ \\
  \hline
\end{tabular}
\end{center}

\medskip

\textbf{Problème 1 --- Est-ce bien une densité ?}
On modélise la durée de vie $X$ (en années) d'un composant par
\[
  f(x)=
  \begin{cases}
    c\,(10-x) & \text{si } 0\le x\le 10,\\
    0 & \text{sinon.}
  \end{cases}
\]
\begin{enumerate}
  \item Déterminer la constante $c$ pour que $f$ soit une densité.
  \item Calculer la probabilité que le composant tombe en panne
  avant $2$ ans.
\end{enumerate}

\atrou[5.5cm]{%
\begin{enumerate}
\item On impose la positivité (ok dès que $c\ge 0$)
et que l'aire vaille $1$ :
\[
  \int_{-\infty}^{+\infty}\! f(x)\,dx
  = \int_0^{10} c(10-x)\,dx
  = c\,\Bigl[10x-\tfrac{x^2}{2}\Bigr]_0^{10}
  = c\cdot 50 = 1
  \;\;\Longrightarrow\;\; c=\tfrac{1}{50}.
\]
\item
\[
  P(X\le 2)
  = \int_0^2 \tfrac{1}{50}(10-x)\,dx
  = \tfrac{1}{50}\Bigl[10x-\tfrac{x^2}{2}\Bigr]_0^2
  = \tfrac{1}{50}\cdot 18 = 0{,}36.
\]
Environ $36\,\%$ des composants tombent en panne
dans les deux premières années.
\end{enumerate}
}

% ====================================================================
\section*{4. Un paradoxe apparent : $P(X=a)=0$}
% ====================================================================

Prenons $a=b$ dans la formule encadrée :
\[
  P(X=a)=P(a\le X\le a)=\int_a^a f(x)\,dx = 0.
\]

\textbf{Comment est-ce possible ?}
Pourtant $X$ prend bien \emph{une} valeur quand on la mesure !

\medskip

\textbf{Intuition.} Il y a une infinité non dénombrable de valeurs
possibles ; si chacune avait une probabilité strictement positive,
la somme exploserait. En continu, la probabilité ne vit pas
\og en un point \fg{}, mais \og sur un intervalle \fg{}.

\medskip

\textbf{Conséquence pratique.}
Pour une variable continue, les inégalités strictes et larges
donnent la même probabilité :
\[
  P(X<a)=P(X\le a),\qquad P(a<X<b)=P(a\le X\le b)
  =\int_a^b f(x)\,dx.
\]

\textbf{Problème 2 --- Lecture d'une densité.}
Soit $X$ de densité $f(x)=\dfrac{3}{4}(1-x^2)$ pour $x\in[-1,1]$,
et $f(x)=0$ sinon.
\begin{enumerate}
  \item Vérifier que $f$ est bien une densité.
  \item Calculer $P(X\ge 0{,}5)$.
  \item Calculer $P(|X|\le 0{,}5)$.
\end{enumerate}
\atrou[6cm]{%
\begin{enumerate}
\item $f\ge 0$ sur $[-1,1]$ (car $1-x^2\ge 0$).
De plus
\[
  \int_{-1}^{1}\tfrac{3}{4}(1-x^2)\,dx
  =\tfrac{3}{4}\Bigl[x-\tfrac{x^3}{3}\Bigr]_{-1}^{1}
  =\tfrac{3}{4}\cdot\tfrac{4}{3}=1.\;\checkmark
\]
\item 
$P(X\ge 0{,}5)=\displaystyle\int_{0{,}5}^{1}\tfrac{3}{4}(1-x^2)\,dx
=\tfrac{3}{4}\Bigl[x-\tfrac{x^3}{3}\Bigr]_{0{,}5}^{1}
=\tfrac{3}{4}\bigl(\tfrac{2}{3}-\tfrac{11}{24}\bigr)
=\tfrac{5}{32}\approx 0{,}156.$

\item Par symétrie ou calcul direct :
$P(|X|\le 0{,}5)=\displaystyle\int_{-0{,}5}^{0{,}5}\tfrac{3}{4}(1-x^2)\,dx
=\tfrac{11}{16}\approx 0{,}688.$
\end{enumerate}
}

% ====================================================================
\section*{5. Fonction de répartition}
% ====================================================================

En discret, on manipulait déjà des \emph{cumuls} :
\(P(X\le 3)=p_1+p_2+p_3\).
On définit l'analogue en continu.

\begin{mdframed}[backgroundcolor=ThemeLight, linecolor=Theme,
  linewidth=1pt, innertopmargin=8pt, innerbottommargin=8pt, nobreak=true]
\textbf{Définition.}
La \textbf{fonction de répartition} de $X$ est
\[
  \boxed{F(x)=P(X\le x).}
\]
Si $X$ admet une densité $f$, alors
\[
  F(x)=\int_{-\infty}^{x} f(t)\,dt
  \qquad\text{et, là où $F$ est dérivable,}\qquad
  f(x)=F'(x).
\]

On note souvent ces quantités avec le nom de la variable aléatoire en indice : 
\[
  F_X(x)=P(X\le x), \quad f_X(x)=F_X'(x).
\]
\end{mdframed}

\textbf{Propriétés}
$F$ est croissante, continue, avec
\[\lim_{x\to-\infty}F(x)=0
\qquad \text{et} \qquad \lim_{x\to+\infty}F(x)=1.\]

\atrou[1.5cm]{%

}
\medskip

\vspace{1em}
\textbf{Problème 3 --- De $f$ à $F$, et calcul de probabilités.}

On prend $f(x)=\tfrac{1}{50}(10-x)$ sur $[0,10]$
(modélisant la durée de vie d'un composant).
\begin{enumerate}
  \item Déterminer $F(x)$ pour tout $x\in\R$.
  \item En déduire $P(3\le X\le 7)$ sans refaire d'intégrale.
\end{enumerate}
\atrou[6cm]{%
\begin{enumerate}
\item Pour $x<0$, $F(x)=0$. Pour $x>10$, $F(x)=1$.
Pour $x\in[0,10]$ :
\[
  F(x)=\int_0^x \tfrac{1}{50}(10-t)\,dt
  =\tfrac{1}{50}\Bigl[10t-\tfrac{t^2}{2}\Bigr]_0^x
  =\tfrac{x}{5}-\tfrac{x^2}{100}
  =\tfrac{x(20-x)}{100}.
\]
\item Le \textbf{réflexe clé} :
\[
  P(a\le X\le b)=F(b)-F(a).
\]
Donc
\(
  P(3\le X\le 7)=F(7)-F(3)
  =\tfrac{7\cdot 13}{100}-\tfrac{3\cdot 17}{100}
  =\tfrac{91-51}{100}=0{,}40.
\)
\end{enumerate}
}

\medskip

\vspace{1em}
\textbf{Problème 4 --- De $F$ à $f$.}

La répartition du temps
d'attente $T$ (en minutes) à un poste de contrôle est donnée par la fonction de répartition:
\[
  F_T(t)=
  \begin{cases}
    0 & \text{si } t<0,\\
    1-\mathrm{e}^{-t/5} & \text{si } t\ge 0.
  \end{cases}\]
\begin{enumerate}
  \item Vérifier que $F_T$ est bien une fonction de répartition.
  \item En déduire la densité $f_T$ de $T$.
  \item Calculer $P(T>10)$.
\end{enumerate}

\atrou[6.5cm]{%
\begin{enumerate}
\item $F_T$ est continue, croissante (sa dérivée
sera positive), $F(t)\to 0$ en $-\infty$ et
$F(t)\to 1$ en $+\infty$. \checkmark

\item Là où $F_T$ est dérivable :
$f(t)=F'(t)=\tfrac{1}{5}\mathrm{e}^{-t/5}$ pour $t\ge 0$,
et $f(t)=0$ pour $t<0$.

\item $P(T>10)=1-F_T(10)=\mathrm{e}^{-2}\approx 0{,}135.$
\end{enumerate}
}

% ====================================================================
\section*{6. Densité d'une fonction de variable aléatoire}
% ====================================================================

Souvent, on a besoin de transformer les variables aléatoires : 

Si on connaît la loi d'une tension $X$, on voudra sans doute
la loi de la puissance $Y=X^2$ ; 

si on connaît le rayon $R$
d'une section voudra la loi de l'aire $A=\pi R^2$.

\medskip

\textbf{La question.}
Si $X$ a pour densité $f_X$, et si $Y=g(X)$, quelle est la densité $f_Y$ ?

\medskip

\textbf{L'idée.}
On ne sait pas calculer directement $f_Y$, mais on sait calculer
la \emph{répartition} $F_Y(y)=P(Y\le y)$ en repassant à $X$.
Puis on dérive.

\begin{mdframed}[backgroundcolor=ThemeLight, linecolor=Theme,
  linewidth=1pt, innertopmargin=8pt, innerbottommargin=8pt, nobreak=true]
\textbf{Méthode (via la fonction de répartition).}
Pour déterminer la densité de $Y=g(X)$ :
\begin{enumerate}
  \item écrire $F_Y(y)=P\bigl(g(X)\le y\bigr)$ ;
  \item \textbf{résoudre} l'inégalité $g(x)\le y$ pour faire apparaître
    un intervalle sur~$x$ ;
  \item exprimer $F_Y$ à l'aide de $F_X$ (ou comme une intégrale de $f_X$) ;
  \item dériver : $f_Y(y)=F_Y'(y)$.
\end{enumerate}
\end{mdframed}

\textbf{Problème 5 --- Transformation affine.}
Soit $X$ une variable continue de densité $f_X$.
On pose $Y=aX+b$ avec $a>0$.
Exprimer $F_Y$ en fonction de $F_X$, puis $f_Y$ en fonction de $f_X$.

\atrou[5cm]{%
\[
  F_Y(y)=P(aX+b\le y)
  =P\!\left(X\le \tfrac{y-b}{a}\right)
  =F_X\!\left(\tfrac{y-b}{a}\right).
\]
On dérive par la règle de la chaîne :
\[
  \boxed{\,f_Y(y)=\tfrac{1}{a}\,f_X\!\left(\tfrac{y-b}{a}\right).\,}
\]
(Dans le cas général $a\ne 0$, on obtient $f_Y(y)=\frac{1}{|a|}f_X\!\bigl(\tfrac{y-b}{a}\bigr)$.)
}

% ====================================================================
\section*{7. Espérance : le centre de gravité, version continue}
% ====================================================================

En CM1 on a vu : \emph{l'espérance, c'est le centre de gravité
de la loi.} Les \og masses \fg{} étaient les probabilités $p_i$
placées aux positions $x_i$ :
\(\,E[X]=\sum_i x_i\, p_i.\)

\medskip

Que devient cette idée si la masse est étalée continûment
le long d'une tige ? 

\atrou[1.5cm]{%
On découpe la tige en petites tranches
de largeur $dx$ autour du point $x$ ; cette tranche porte
une \og masse \fg{} $f(x)\,dx$ (aire de la petite bande
de l'histogramme). Le centre de gravité est alors :
\[
  \underbrace{\sum_i x_i\,p_i}_{\text{discret}}
  \quad\longrightarrow\quad
  \underbrace{\int_{-\infty}^{+\infty} x\,f(x)\,dx}_{\text{continu}}.
\]
}

\begin{mdframed}[backgroundcolor=ThemeLight, linecolor=Theme,
  linewidth=1pt, innertopmargin=8pt, innerbottommargin=8pt, nobreak=true]
\textbf{Définition.}
Si $X$ admet une densité $f$, son \textbf{espérance} est
\[
  \boxed{\,E[X]=\int_{-\infty}^{+\infty} x\,f(x)\,dx\,}
\]
(lorsque cette intégrale est absolument convergente).
\end{mdframed}

\textbf{Problème 6 --- Calcul d'espérances.}
\begin{enumerate}
  \item Durée de vie $X$ de densité
  $f(x)=\tfrac{1}{50}(10-x)$ sur $[0,10]$ : calculer $E[X]$.
  Interpréter.
  \item Temps d'attente $T$ de densité
  $f(t)=\tfrac{1}{5}\mathrm{e}^{-t/5}$ sur $[0,+\infty[$ :
  calculer $E[T]$ (on admettra
  $\int_0^{+\infty} t\,\mathrm{e}^{-t/\lambda}\,dt = \lambda^2$).
\end{enumerate}

\atrou[6cm]{%
\begin{enumerate}
\item
\[
  E[X]=\int_0^{10} x\cdot\tfrac{1}{50}(10-x)\,dx
  =\tfrac{1}{50}\Bigl[5x^2-\tfrac{x^3}{3}\Bigr]_0^{10}
  =\tfrac{1}{50}\cdot\tfrac{500}{3}=\tfrac{10}{3}\approx 3{,}33\text{ ans}.
\]
En moyenne, un composant dure $3$ ans et $4$ mois.
La densité décroît, donc les petites valeurs sont plus
probables : le centre de gravité est bien dans la moitié
gauche de l'intervalle.

\item
\(\displaystyle
  E[T]=\int_0^{+\infty}\! t\cdot\tfrac{1}{5}\mathrm{e}^{-t/5}\,dt
  =\tfrac{1}{5}\cdot 25 = 5\text{ min}.
\)
\end{enumerate}
}

\medskip

\textbf{Théorème (de transfert).}
Pour calculer l'espérance d'une \emph{fonction} de $X$, pas besoin
de trouver la densité de $g(X)$ : on intègre directement,
\[
  \boxed{\,E[g(X)]=\int_{-\infty}^{+\infty} g(x)\,f(x)\,dx.\,}
\]
(Version continue de la formule $E[g(X)]=\sum_i g(x_i)\,p_i$
du cas discret.)

% ====================================================================
\section*{8. Applications du théorème de transfert}
% ====================================================================

Le théorème de transfert a deux usages pratiques majeurs.

\medskip

\textbf{1. Calculer directement $E[g(X)]$} sans déterminer la densité
de $g(X)$ : on intègre $g(x)\,f_X(x)$. On s'en servira notamment
pour calculer $E[X^2]$, qui intervient dans la variance (§9).

\medskip

\textbf{2. Retrouver la densité de $Y=g(X)$} sans passer par la
fonction de répartition. C'est la \og méthode de la fonction muette \fg{}.
L'idée : pour toute fonction continue bornée $h$, le transfert donne
\[
  E[h(Y)] = E\bigl[h(g(X))\bigr]
  = \int_{-\infty}^{+\infty} h(g(x))\,f_X(x)\,dx.
\]
Si on arrive à réécrire cette intégrale, par un changement de variable
$y=g(x)$, sous la forme $\int h(y)\,\varphi(y)\,dy$, alors par
comparaison avec $E[h(Y)] = \int h(y)\,f_Y(y)\,dy$ (vraie pour
\emph{tout} $h$), on identifie directement $f_Y=\varphi$.

\begin{mdframed}[backgroundcolor=ThemeLight, linecolor=Theme,
  linewidth=1pt, innertopmargin=8pt, innerbottommargin=8pt, nobreak=true]
\textbf{Méthode (via le transfert, dite \og de la fonction muette \fg).}
Pour déterminer la densité de $Y=g(X)$ :
\begin{enumerate}
  \item pour une fonction $h$ continue bornée \emph{quelconque},
    écrire $E[h(Y)]=\int h(g(x))\,f_X(x)\,dx$ grâce au transfert ;
  \item effectuer le changement de variable $y=g(x)$ dans l'intégrale
    pour obtenir une écriture de la forme $\int h(y)\,\varphi(y)\,dy$ ;
  \item identifier avec $\int h(y)\,f_Y(y)\,dy$ : l'égalité étant
    valable pour tout~$h$, on a $f_Y(y)=\varphi(y)$.
\end{enumerate}
\end{mdframed}

\textbf{Problème 7 --- Passage au carré.}
Soit $X$ de densité $f_X(x)=1$ pour $x\in[0,1]$, et $0$ sinon
(chaque valeur de $[0,1]$ est \og également plausible \fg{}).
On pose $Y=X^2$.
\begin{enumerate}
  \item Dans quel intervalle $Y$ prend-elle ses valeurs ?
  \item Déterminer $F_Y(y)$ pour $y\in[0,1]$, puis $f_Y$ sur~$\R$
    (méthode du §6).
  \item Calculer $E[Y]$ de deux façons : directement avec $f_Y$, puis
    \emph{via le théorème de transfert}. Comparer.
\end{enumerate}

\atrou[7.5cm]{%
\begin{enumerate}
\item $X\in[0,1]$ donc $Y=X^2\in[0,1]$.

\item Pour $y\in[0,1]$, comme $X\ge 0$ :
\[
  F_Y(y)=P(X^2\le y)=P(0\le X\le\sqrt{y})
  =\int_0^{\sqrt{y}}\! 1\,dx=\sqrt{y}.
\]
Pour $y<0$, $F_Y(y)=0$ ; pour $y>1$, $F_Y(y)=1$.
En dérivant :
\[
  f_Y(y)=
  \begin{cases}
    \dfrac{1}{2\sqrt{y}} & \text{si } y\in\,]0,1],\\
    0 & \text{sinon.}
  \end{cases}
\]
(Noter que $f_Y$ tend vers $+\infty$ en $0$ : les petites valeurs
de $Y$ sont \og concentrées \fg{}.)

\item Via $f_Y$ :
$\;E[Y]=\int_0^1 y\cdot\tfrac{1}{2\sqrt{y}}\,dy
=\tfrac{1}{2}\int_0^1\!\sqrt{y}\,dy=\tfrac{1}{2}\cdot\tfrac{2}{3}=\tfrac{1}{3}.$

Via le transfert, sans passer par $f_Y$ :
$\;E[Y]=E[X^2]=\int_0^1 x^2\cdot 1\,dx=\tfrac{1}{3}.\;\checkmark$
\end{enumerate}
}

\medskip

\textbf{Problème 8 --- La méthode de la fonction muette.}
Soit $X$ de densité $f_X(x)=2x$ pour $x\in[0,1]$, $0$ sinon.
On pose $Y=\sqrt{X}$, et on veut la densité de $Y$ par la méthode
de la fonction muette.
\begin{enumerate}
  \item Soit $h$ une fonction continue bornée. Écrire $E[h(Y)]$ comme
    une intégrale en $x$ grâce au théorème de transfert.
  \item Effectuer le changement de variable $y=\sqrt{x}$ (soit $x=y^2$)
    et donner la nouvelle intégrale.
  \item Par identification, en déduire la densité $f_Y$.
\end{enumerate}

\atrou[7cm]{%
\begin{enumerate}
\item Avec $g(x)=\sqrt{x}$, le transfert donne
\[
  E[h(Y)] = E\bigl[h(\sqrt{X})\bigr]
  = \int_0^1 h(\sqrt{x})\cdot 2x\,dx.
\]
\item On pose $y=\sqrt{x}$, donc $x=y^2$ et $dx=2y\,dy$.
Quand $x$ parcourt $[0,1]$, $y$ parcourt aussi $[0,1]$. Donc
\[
  E[h(Y)] = \int_0^1 h(y)\cdot 2y^2\cdot 2y\,dy
  = \int_0^1 h(y)\cdot 4y^3\,dy.
\]
\item Cette égalité est valable pour \emph{toute} fonction $h$, donc
par identification avec $E[h(Y)]=\int_\R h(y)\,f_Y(y)\,dy$ :
\[
  f_Y(y)=
  \begin{cases}
    4y^3 & \text{si } y\in[0,1],\\
    0 & \text{sinon.}
  \end{cases}
\]
\textit{Vérification :} $\int_0^1 4y^3\,dy=1$. \checkmark
\end{enumerate}

\textbf{Bilan.} On obtient la densité sans avoir à déterminer
l'intervalle $\{g(X)\le y\}$ ni à dériver. Cette méthode est souvent
plus rapide quand $g$ est une bijection régulière ; en revanche, quand
$g$ n'est pas monotone (par exemple $g(x)=x^2$ sur un intervalle
contenant $0$), il faut découper le changement de variable, et la
méthode via $F_Y$ du §6 peut être plus simple.
}

% ====================================================================
\section*{9. Variance : le moment d'inertie continu}
% ====================================================================

Même musique qu'en CM1 : une fois qu'on connaît le centre
de gravité $E[X]$, on mesure la \textbf{dispersion} autour
de lui par un moment d'inertie.

\begin{mdframed}[backgroundcolor=ThemeLight, linecolor=Theme,
  linewidth=1pt, innertopmargin=8pt, innerbottommargin=8pt, nobreak=true]
\textbf{Définition.}
Si $X$ a une densité $f$ et une espérance $E[X]$, sa \textbf{variance}
est
\[
  \boxed{\text{Var}(X)=E\bigl[(X-E[X])^2\bigr]
  =\int_{-\infty}^{+\infty}(x-E[X])^2\,f(x)\,dx.}
\]
L'\textbf{écart-type} est $\sigma(X)=\sqrt{\text{Var}(X)}$
(même unité que $X$).
\end{mdframed}

\textbf{La formule pratique (Huygens) reste vraie.}
En développant le carré, comme en CM1 :
\[
  \boxed{\text{Var}(X)=E[X^2]-E[X]^2}
  \qquad\text{avec}\qquad
  E[X^2]=\int_{-\infty}^{+\infty} x^2\,f(x)\,dx.
\]

\medskip

\textbf{Problème 9 --- Variance de la durée de vie.}
Reprenons $X$ avec $f(x)=\tfrac{1}{50}(10-x)$ sur $[0,10]$,
$E[X]=10/3$.

\atrou[5.5cm]{%
\[
  E[X^2]
  =\int_0^{10} x^2\cdot\tfrac{1}{50}(10-x)\,dx
  =\tfrac{1}{50}\Bigl[\tfrac{10x^3}{3}-\tfrac{x^4}{4}\Bigr]_0^{10}
  =\tfrac{1}{50}\cdot\tfrac{2500}{6}
  =\tfrac{50}{3}.
\]
\[
  \text{Var}(X)=\tfrac{50}{3}-\Bigl(\tfrac{10}{3}\Bigr)^{\!2}
  =\tfrac{50}{3}-\tfrac{100}{9}
  =\tfrac{150-100}{9}
  =\tfrac{50}{9}\approx 5{,}56.
\]
$\sigma(X)\approx 2{,}36$ ans.
}

\textbf{Problème 10 --- Un test de cohérence.}
Pour la densité \og parabolique \fg{}
$f(x)=\tfrac{3}{4}(1-x^2)$ sur $[-1,1]$ (problème~2),
expliquer \emph{sans calcul} pourquoi $E[X]=0$,
puis vérifier en calculant. Calculer ensuite $\text{Var}(X)$.

\atrou[5.5cm]{%
La densité est \textbf{paire} ($f(-x)=f(x)$), donc son
centre de gravité est $0$ par symétrie. Calcul :
$E[X]=\int_{-1}^{1} x\cdot\tfrac{3}{4}(1-x^2)\,dx=0$
(intégrale d'une fonction impaire sur $[-1,1]$). \checkmark
\[
  E[X^2]=\int_{-1}^{1} x^2\cdot\tfrac{3}{4}(1-x^2)\,dx
  =\tfrac{3}{4}\Bigl[\tfrac{x^3}{3}-\tfrac{x^5}{5}\Bigr]_{-1}^{1}
  =\tfrac{3}{4}\cdot\tfrac{4}{15}=\tfrac{1}{5}.
\]
Donc $\text{Var}(X)=\tfrac{1}{5}-0^2=0{,}20$,
$\sigma(X)\approx 0{,}447$.
}

% ====================================================================
\section*{10. Propriétés : ce qui ne change pas par rapport au discret}
% ====================================================================

Les propriétés vues au CM1 restent valables : elles découlent
de la définition de $E$ et $\text{Var}$, pas du type
(discret ou continu) de la variable.

\medskip

\textbf{Linéarité de l'espérance.}
Pour $a,b\in\R$,
\[
  E[aX+b]=a\,E[X]+b,
  \qquad
  E[X+Y]=E[X]+E[Y].
\]

\textbf{Variance et transformation affine.}
\[
  \text{Var}(aX+b)=a^2\,\text{Var}(X),
  \qquad
  \sigma(aX+b)=|a|\,\sigma(X).
\]

\textbf{Problème 11 --- Changement d'unité.}

\textit{(Ce problème illustre aussi la méthode du §6, avec une
transformation affine.)}
La longueur $X$ d'une pièce est mesurée en cm, avec
$E[X]=20$ cm et $\sigma(X)=0{,}5$ cm.
On définit l'\textbf{écart au nominal} $Y=X-20$ (en mm,
sachant que $1$ cm $=10$ mm, donc $Y=10(X-20)$ en mm).
Calculer $E[Y]$ et $\sigma(Y)$.

\atrou[3.5cm]{%
$E[Y]=E[10(X-20)]=10\,(E[X]-20)=10\cdot 0=0$ mm.

$\sigma(Y)=|10|\cdot\sigma(X)=10\cdot 0{,}5=5$ mm.
L'écart au nominal est en moyenne nul, avec un écart-type
de $5$ mm.
}

% ====================================================================
\section*{11. Bilan}
% ====================================================================

\begin{center}
\renewcommand{\arraystretch}{1.6}
\begin{tabular}{|l|c|c|}
  \hline
  & \textbf{Discret (CM1)} & \textbf{Continu (CM2)} \\
  \hline
  Objet associé à $X$ & probabilités $p_i=P(X=x_i)$ & densité $f(x)$ \\
  \hline
  Normalisation & $\sum_i p_i = 1$
    & $\displaystyle\int_{-\infty}^{+\infty} f(x)\,dx=1$ \\
  \hline
  Probabilité d'un événement & $\sum_{x_i\in A} p_i$
    & $\displaystyle\int_A f(x)\,dx$ \\
  \hline
  Répartition & $F(x)=\sum_{x_i\le x}p_i$
    & $F(x)=\int_{-\infty}^{x} f(t)\,dt$ \\
  \hline
  Espérance & $\sum_i x_i\,p_i$
    & $\displaystyle\int_{-\infty}^{+\infty} x\,f(x)\,dx$ \\
  \hline
  Variance & $E[X^2]-E[X]^2$
    & $E[X^2]-E[X]^2$ \\
  \hline
\end{tabular}
\end{center}

\medskip

\atrou[6cm]{%
\begin{enumerate}
  \item Une variable continue se manipule avec une \textbf{densité}~$f$ :
    positive, d'aire totale $1$, et
    $P(a\le X\le b)=\int_a^b f(x)\,dx$.
  \item En continu, $P(X=a)=0$ : seules les probabilités
    \textbf{d'intervalles} ont un sens non trivial.
    Les bornes strictes ou larges donnent la même valeur.
  \item La \textbf{répartition} $F(x)=P(X\le x)$ s'obtient en intégrant $f$,
    et inversement $f=F'$ ; la formule clé reste
    $P(a\le X\le b)=F(b)-F(a)$.
  \item Pour la \textbf{densité de $Y=g(X)$}, on passe par la
    répartition : $F_Y(y)=P(g(X)\le y)$, qu'on exprime en fonction
    de $F_X$, puis on dérive.
  \item \textbf{Espérance} et \textbf{variance} sont définies comme en discret,
    en remplaçant \og $\sum p_i\cdot\ldots$ \fg{} par
    \og $\int f(x)\cdot\ldots\,dx$ \fg{}.
    Les propriétés (linéarité, Huygens, $\text{Var}(aX+b)=a^2\text{Var}(X)$)
    sont les mêmes.
  \item \textbf{Au CM3}, nous découvrirons plusieurs lois continues
    \emph{classiques} (uniforme, exponentielle, \dots) qui jouent le rôle
    des lois de Bernoulli, binomiale, géométrique, Poisson du CM1.
\end{enumerate}
}

\end{document}
